% 2013 年辽宁卷理科第 5 题：成绩频率分布直方图。
% 来源 PDF：普通高考/2013/2013辽宁理.pdf，第 1 页（印面 821）左栏。
% 分组 [20,40)、[40,60)、[60,80)、[80,100] 的柱高依次为
% 0.005、0.010、0.020、0.015（300 dpi 原页核对）。
\begin{tikzpicture}[font=\footnotesize, >=Stealth, line width=0.5pt,
  x=1.15cm, y=160cm]
  % 坐标轴：横轴以 20 分为一个单位
  \draw[->] (0,0) -- (5.75,0) node[right, yshift=-2pt] {成绩/分};
  \draw[->] (0,0) -- (0,0.0238);
  \node at (0,0.0295) {$\dfrac{\text{频率}}{\text{组距}}$};
  % 横轴刻度
  \foreach \x/\lab in {0/0,1/20,2/40,3/60,4/80,5/100}
    \draw (\x,0) -- (\x,-0.0006) node[below] {$\lab$};
  % 纵轴刻度与虚线
  \foreach \y/\lab in {0.005/0.005,0.010/0.010,0.015/0.015,0.020/0.020} {
    \draw (0,\y) -- (-0.08,\y) node[left] {$\lab$};
    \draw[dashed] (0,\y) -- (5.05,\y);
  }
  % 四个矩形位于 [20,40)、[40,60)、[60,80)、[80,100]，
  % 高度依次为 0.005、0.010、0.020、0.015
  \foreach \i/\h in {1/0.005,2/0.010,3/0.020,4/0.015} {
    \draw[fill=white] (\i,0) rectangle (\i+1,\h);
  }
\end{tikzpicture}
